% ---------------------------------------------------------------------------
% What is left out and why: (a) the three gg -> b bbar diagrams that
% HardQCD:gg2bbbar = off removes, (b) the two b bbar -> b bbar diagrams of
% which PYTHIA process 124 keeps only the first, (c) the five-flavour picture
% of a b quark inside the proton (g -> b bbar in the evolution) against the
% four-flavour one, where the same splitting sits in the matrix element.
% ---------------------------------------------------------------------------
\begin{tikzpicture}[x=1cm,y=1cm,
  every node/.style={font=\footnotesize},
  fermion/.style={thick,postaction={decorate},decoration={markings,mark=at position 0.55 with {\arrow[scale=1.0]{Stealth}}}},
  antifermion/.style={thick,postaction={decorate},decoration={markings,mark=at position 0.45 with {\arrowreversed[scale=1.0]{Stealth}}}},
  plain/.style={thick},
  gluon/.style={thick,decorate,decoration={coil,aspect=0,segment length=4pt,amplitude=1.8pt}},
  vertex/.style={fill=black,circle,inner sep=1.4pt},
  lab/.style={font=\scriptsize,align=center},
  head/.style={font=\small\bfseries},
  off/.style={draw=red!60!black,dashed,rounded corners=3pt,thick},
  on/.style={draw=green!50!black,rounded corners=3pt,thick}]
% ---------------- (a) gg -> b bbar: switched off ---------------------------------------------------
\node[head,anchor=west] at (-0.3,2.6) {(a) $gg\to\bbbar$: three diagrams, switched off in our run};
\begin{scope}[xshift=0cm]
\draw[gluon] (0,1.2) -- (1.2,0.6); \draw[gluon] (0,-1.2) -- (1.2,-0.6);
\draw[gluon] (1.2,0.6) -- (1.2,-0.6);
\draw[fermion] (1.2,0.6) -- (2.4,1.2) node[right,lab] {$b$}; \draw[antifermion] (1.2,-0.6) -- (2.4,-1.2) node[right,lab] {$\bar b$};
\node[vertex] at (1.2,0.6) {}; \node[vertex] at (1.2,-0.6) {};
\node[lab,below] at (1.2,-1.4) {$t$-channel $b$};
\end{scope}
\begin{scope}[xshift=3.6cm]
\draw[gluon] (0,1.2) -- (1.2,0.6); \draw[gluon] (0,-1.2) -- (1.2,-0.6);
\draw[gluon] (1.2,0.6) -- (1.2,-0.6);
\draw[fermion] (1.2,0.6) -- (2.4,-1.2) node[right,lab] {$b$}; \draw[antifermion] (1.2,-0.6) -- (2.4,1.2) node[right,lab] {$\bar b$};
\node[vertex] at (1.2,0.6) {}; \node[vertex] at (1.2,-0.6) {};
\node[lab,below] at (1.2,-1.4) {$u$-channel $b$};
\end{scope}
\begin{scope}[xshift=7.2cm]
\draw[gluon] (0,1.2) -- (0.8,0); \draw[gluon] (0,-1.2) -- (0.8,0);
\draw[gluon] (0.8,0) -- (1.8,0);
\draw[fermion] (1.8,0) -- (2.8,1.2) node[right,lab] {$b$}; \draw[antifermion] (1.8,0) -- (2.8,-1.2) node[right,lab] {$\bar b$};
\node[vertex] at (0.8,0) {}; \node[vertex] at (1.8,0) {};
\node[lab,below] at (1.4,-1.4) {$s$-channel $g$,\\three-gluon vertex};
\end{scope}
\draw[off] (-0.4,-2.4) rectangle (10.4,1.6);
\node[lab,text=red!60!black,anchor=north] at (5.0,-2.5) {\code{HardQCD:gg2bbbar = off}: dominant in an inclusive $\bbbar$ sample, and not a one-line formula};
% ---------------- (b) b bbar -> b bbar: only the s-channel is used -----------------------------------
\begin{scope}[xshift=11.6cm]
\node[head,anchor=west] at (-0.3,2.6) {(b) the incoming $b$: $b\bar b\to b\bar b$};
\begin{scope}
\draw[fermion] (0,1.2) -- (0.8,0); \draw[antifermion] (0,-1.2) -- (0.8,0);
\draw[gluon] (0.8,0) -- (1.8,0);
\draw[fermion] (1.8,0) -- (2.6,1.2); \draw[antifermion] (1.8,0) -- (2.6,-1.2);
\node[vertex] at (0.8,0) {}; \node[vertex] at (1.8,0) {};
\node[lab,below] at (1.3,-1.4) {$s$-channel: used\\(process 124,\\all five flavours)};
\draw[on] (-0.3,-2.4) rectangle (2.9,1.6);
\end{scope}
\begin{scope}[xshift=4.2cm]
\draw[fermion] (0,1.2) -- (1.0,0.6) -- (2.2,1.2); \draw[antifermion] (0,-1.2) -- (1.0,-0.6) -- (2.2,-1.2);
\draw[gluon] (1.0,0.6) -- (1.0,-0.6);
\node[vertex] at (1.0,0.6) {}; \node[vertex] at (1.0,-0.6) {};
\node[lab,below] at (1.1,-1.4) {$t$-channel: ignored\\(part of\\the $1.6\%$)};
\draw[off] (-0.3,-2.4) rectangle (2.5,1.6);
\end{scope}
\end{scope}
% ---------------- (c) five against four flavours -------------------------------------------------------
\begin{scope}[yshift=-6.2cm]
\node[head,anchor=west] at (-0.3,2.3) {(c) where the $b$ of the proton comes from: five-flavour against four-flavour scheme};
\begin{scope}
\node[draw,ellipse,thick,fill=gray!10,minimum width=16mm,minimum height=8mm] (P) at (0.6,0) {$p$};
\draw[gluon] (P.east) -- (2.4,0);
\draw[fermion] (2.4,0) -- (3.6,0.9) node[right,lab] {$b$ enters $\hat\sigma$ (massless)};
\draw[antifermion] (2.4,0) -- (3.6,-0.9) node[right,lab] {$\bar b$ stays in the remnant};
\node[vertex] at (2.4,0) {};
\node[lab,text=black!65,text width=7.4cm,anchor=north west] at (-0.3,-1.3) {5FS (\PYTHIA, CMS, the scripts): the $g\to\bbbar$ splitting is inside the PDF $f_b(x,\muF^2)$, with its $\ln(\muF^2/\mb^2)$ resummed by DGLAP; $f_b=0$ below $\mb$. Incoming-$b$ share here: $1.6\%$.};
\end{scope}
\begin{scope}[xshift=11.6cm]
\node[draw,ellipse,thick,fill=gray!10,minimum width=16mm,minimum height=8mm] (P) at (0.6,0) {$p$};
\draw[gluon] (P.east) -- (2.4,0);
\draw[fermion] (2.4,0) -- (3.6,0.9) node[right,lab] {$b$ in the matrix element};
\draw[antifermion] (2.4,0) -- (3.6,-0.9) node[right,lab] {$\bar b$ in the final state};
\node[vertex] at (2.4,0) {};
\node[draw,circle,thick,fill=orange!25,minimum size=7mm] at (4.9,0.0) {$\hat\sigma$};
\node[lab,text=black!65,text width=7.4cm,anchor=north west] at (-0.3,-1.3) {4FS (\MG's default proton): no $b$ in the proton; the same splitting is an explicit $2\to3$ (or higher) matrix element with the $b$ mass kept. Equal to 5FS once higher orders are included; $1.6\%$ lower at this order.};
\end{scope}
\end{scope}
\end{tikzpicture}
